\section{NP-hardness of approximation}
\label{sec:np-hardness}

We prove NP-hardness below $18/17$ for treedepth, below $12/11$
for treewidth, and below $63/62$ for total elimination depth.
All three reductions use the normalized bounded-occurrence parity source
of Theorem~\ref{thm:xor-source}. Treedepth and treewidth share one
tripartite construction. A second tripartite encoding gives the
total-depth gap in Section~\ref{subsec:np-total-depth}.

\subsection{Treedepth}
\label{subsec:np-treedepth}

We encode parity equations in a tripartite graph, then convert its
vertex-cover number into a dense treedepth gap.
Corollary~\ref{cor:td-bounded-degree} gives the bounded-degree
hardness consequence after sparsification.

The truth blocks and assignment witnesses follow Bonnet, Neuen, and
Soko\l owski~\cite[Section~4 and Lemma~10]{BNS}, building on Bonnet's
treewidth construction~\cite[Section~3.1]{BonnetTreewidth}. The anchor graph
is the one in~\cite[Lemma~6]{BNS}. We modify the witness gadgets
for parity equations and use the anchor construction through the
corollary below. Appendix~\ref{subsec:np-provenance} records which
ingredients are inherited, and Appendix~\ref{app:anchor-proof} gives
the corollary's derivation.

\subsubsection{The source gap and normalization}

We use H{\aa}stad's Max-3-LIN-$2$
gap~\cite{Hastad1997}, in the distinct-variable,
bounded-occurrence form stated by Hazan, Safra, and Schwartz
\cite[Theorem~10]{HazanSafraSchwartz2003}.
Appendix~\ref{app:xor-normalization} adds right-hand-side balance and
divisibility of occurrence counts to obtain the following statement.

\begin{theorem}[Normalized bounded-occurrence Max-3-LIN-2 gap]\label{thm:xor-source}
For every fixed $0<\epsilon<1/4$, there is a constant
$b_\epsilon\geq6$ depending only on $\epsilon$ such that the following
promise problem is NP-hard. The input is an unweighted system $F$ with
finite Boolean variable set $\mathcal X(F)$ and $m\geq1$ equations,
indexed by $j\in\{1,\ldots,m\}$, of the form $x+y+z=a\pmod2$,
where $a\in\{0,1\}$ and $x,y,z\in\mathcal X(F)$ are distinct.
Equations may be repeated.
For $x\in\mathcal X(F)$, write $d_x$ for
the number of equations containing $x$. Write $a_j$ for the right-hand
side of equation $j$, and
$\OPT(F)$ for the maximum number simultaneously satisfied,
counting repetitions. The instances satisfy
\begin{equation}\label{eq:xor-normalization}
\begin{gathered}
 d_x\in6\mathbb Z_{>0},\qquad d_x\leq b_\epsilon,\qquad
 |\{d_x:x\in\mathcal X(F)\}|\leq2,\qquad
 \sum_{x\in\mathcal X(F)}d_x=3m,\\
 |\{j\in\{1,\ldots,m\}:a_j=0\}|
 =|\{j\in\{1,\ldots,m\}:a_j=1\}|=m/2.
\end{gathered}
\end{equation}
\textbf{YES:} $\OPT(F)\geq(1-\epsilon)m$;
\textbf{NO:} $\OPT(F)\leq(1/2+\epsilon)m$.
\end{theorem}

For the normalization, let $F_0$ be the initial source system, with
$m_0$ equations. Let $F_0'$ be a copy on fresh variables with every
right-hand side complemented,
and put $F_1=F_0\mathbin{\dot\cup}F_0'$. Lemmas~\ref{lem:parity-balance} and~\ref{lem:parity-repetition}
justify the two operations and their effect on the optimum:
\[
 F_0\xrightarrow{\text{complemented copy}}F_1
 \xrightarrow{\text{sixfold repetition}}F,\qquad (m,\OPT(F))
 \overset{\text{Lems.~\ref{lem:parity-balance}, \ref{lem:parity-repetition}}}{=}
 12(m_0,\OPT(F_0)).
\]
Thus the satisfied fraction is unchanged, and the maximum occurrence
count is multiplied by six.
We retain the correspondence between each variable of $F_0$ and its
copy, and between the repeated equations in the two halves. This
recorded pairing will provide a balanced optimal independent set in
the treewidth reduction. The restriction to at most two occurrence
counts is inherited from the source theorem; the treedepth and
treewidth reductions use only the other normalization conditions.
We henceforth assume~\eqref{eq:xor-normalization}. Divisibility makes
the truth-block sizes integral, and right-hand-side balance makes the
three graph parts equal.

\subsubsection{A tripartite graph with equal parts}

We construct a tripartite graph $G_F$ with independent parts
$A\mathbin{\dot\cup}B\mathbin{\dot\cup}C$. Its vertices have two
roles: vertices in \emph{truth blocks} encode values of individual
variables, and \emph{assignment witnesses} represent satisfying
assignments to the three variables of an equation. All blocks are
pairwise disjoint, and every witness is a new vertex.
For every $x\in\mathcal X(F)$, create independent truth blocks
$T_x^0\subseteq A$ and $T_x^1\subseteq B$. Their size and the total
size of one block per variable are
\begin{equation}\label{eq:xor-truth-budget}
 |T_x^0|=|T_x^1|=5d_x/6\in\mathbb Z_{>0},\qquad \frac56\sum_{x\in\mathcal X(F)}d_x
 \overset{\eqref{eq:xor-normalization}}{=}\frac{5m}{2}.
\end{equation}
Integrality follows from~\eqref{eq:xor-normalization}. Put all edges
between these two blocks. For $a\in\{0,1\}$, selecting $T_x^a$ in an
independent set will encode the assignment $x=a$.

For each equation $j\in\{1,\ldots,m\}$, create four witnesses, one for
each of its satisfying assignments. If its right-hand side is zero, call the
witness for assignment $000$ special and place it in $A$. If its
right-hand side is one, call the witness for assignment $111$ special
and place it in $B$. Place the other three witnesses in $C$. Join the
special witness to the other three witnesses of its equation, forming
a star. We call these four vertices and their star edges the
\emph{equation gadget}. Its star permits at most three witnesses
in an independent set, even when no truth block is selected.

For every witness representing a satisfying assignment $\sigma$ and
every variable $x$ of its equation, join it to the entire truth block
$T_x^{1-\sigma(x)}$. Add no other edges, so witnesses from distinct
equations are nonadjacent.

Every listed edge respects the tripartition. A special $000$ witness
lies in $A$ and its truth-block neighbors all lie in $B$. A special
$111$ witness lies in $B$ and its truth-block neighbors all lie in
$A$. The three remaining witnesses lie in $C$, and their truth-block
neighbors lie in $A\cup B$. Star edges join the special witness to
$C$. Thus $G_F$ is simple, unweighted, and tripartite.

The truth vertices number $5m$, and there are $4m$ witnesses, so
$G_F$ has $9m$ vertices. Equation~\eqref{eq:xor-normalization}
gives $m/2$ special witnesses in each of $A$ and $B$, and three
witnesses per equation in $C$. Hence
\begin{equation}\label{eq:xor-equal-parts}
  |A|=|B|=|C|
  \overset{\eqref{eq:xor-truth-budget},\,\eqref{eq:xor-normalization}}{=}3m.
\end{equation}

The source graph $G_F$ already has bounded maximum degree.
For $v\in T_x^a$, its opposite truth block and witness neighbors give
\[
 \deg_{G_F}(v)=\frac{5d_x}{6}+2d_x
 =\frac{17d_x}{6}
 \overset{\eqref{eq:xor-normalization}}{\leq}3b_\epsilon.
\]
For a witness $w$ of an equation on $x,y,z$, its three truth-block
neighborhoods and at most three star neighbors give
\[
 \deg_{G_F}(w)\leq\frac56(d_x+d_y+d_z)+3
 \overset{\eqref{eq:xor-normalization}}{\leq}\frac52b_\epsilon+3
 \overset{\text{Thm.~\ref{thm:xor-source}}}{\leq}3b_\epsilon,
\]
where the last inequality uses $b_\epsilon\geq6$.

\subsubsection{The independent-set optimum}

\begin{lemma}[Random completion of parity equations]
\label{lem:partial-parity-completion}
Fix any partial assignment to the variables of a system of parity
equations on three distinct variables. Complete it by choosing each
unassigned variable independently and uniformly from $\{0,1\}$.
Every equation with an unassigned variable is satisfied with
probability $1/2$; a fully assigned equation keeps its truth value.
If $s$ fully assigned equations are satisfied and $u$ equations
are unfinished, the expected number satisfied is $s+u/2$, and
some completion satisfies at least this many equations.
\end{lemma}

\begin{proof}
For an unfinished equation, condition on all but one of its
unassigned variables. Its remaining fair bit determines its parity,
so either outcome has probability $1/2$. Sum these probabilities
and the fixed truth values over all equations. Linearity of
expectation does not require independence between equations, and
some outcome attains at least the expectation.
\end{proof}

For an independent set $I$ containing every truth block that it
meets, the selected block $T_x^a$ fixes $x=a$. Variables with neither
block selected remain unassigned. Complete this partial assignment by
choosing each unassigned variable independently and uniformly from
$\{0,1\}$. Denote the resulting random assignment by $\sigma$.
For $j\in\{1,\ldots,m\}$, let $r_j$ count its unassigned variables
and let $w_j$ count its witness vertices in $I$. Define
$p_j=\Pr[\text{equation }j\text{ is satisfied by }\sigma]$.

\begin{lemma}[Witnesses under a partial assignment]
\label{lem:xor-witness-excess}
With this notation, $w_j-p_j\le5r_j/6$ for every $j\in\{1,\ldots,m\}$.
\end{lemma}

\begin{proof}
Every selected witness agrees with the fixed variables, since
otherwise it has a neighbor in a selected truth block.
If $r_j=0$, the equation is either satisfied with one consistent
witness or unsatisfied with none, so $w_j\leq p_j$.
If $r_j\geq1$, Lemma~\ref{lem:partial-parity-completion} gives
$p_j=1/2$. The $2^{r_j}$ completions are equally likely, so exactly
$2^{r_j-1}$ witnesses are consistent. The star permits at most three. Thus, for $r_j\in\{1,2,3\}$,
\begin{equation}\label{eq:xor-local-witness-count}
 w_j\leq\min\{2^{r_j-1},3\}=r_j,\qquad p_j=\frac12,\qquad w_j-p_j
 \overset{\eqref{eq:xor-local-witness-count}}{\leq}
 r_j-\frac12\overset{r_j\leq3}{\leq}\frac56r_j.
\end{equation}
\end{proof}

\begin{lemma}\label{lem:xor-cover}
The graph $G_F$ satisfies
\[
 \Delta(G_F)\leq3b_\epsilon,\qquad
 \alphaI(G_F)=5m/2+\OPT(F),\qquad
 \vc(G_F)=13m/2-\OPT(F).
\]
\end{lemma}


\begin{proof}
Start with any independent set $I$, and fill every truth block that
it meets. This preserves independence because a block has no internal
edges and its vertices have identical neighborhoods. It suffices to
bound the filled set, which we continue to call $I$. The two blocks
of a variable are joined completely, so at most one is selected.
Let $U\subseteq\mathcal X(F)$ be the set of unassigned variables.
Relative to one truth block per variable, the missing truth vertices
number $\frac56\sum_{x\in U}d_x$.
For the random completion above, counting those variables' occurrences
and applying linearity of expectation give
\begin{equation}\label{eq:xor-completion-counts}
 \sum_{j=1}^{m}r_j=\sum_{x\in U}d_x,\qquad
 \sum_{j=1}^{m}p_j\leq\OPT(F).
\end{equation}
The second bound holds because every completion satisfies at most
$\OPT(F)$ equations.
Lemma~\ref{lem:xor-witness-excess} therefore gives
\[
 |I|\overset{\eqref{eq:xor-truth-budget},\,\eqref{eq:xor-completion-counts}}{=}
       \frac{5m}{2}+\sum_{j=1}^{m}\left(w_j-\frac56r_j\right)
 \overset{\text{Lem.~\ref{lem:xor-witness-excess}}}{\le}
       \frac{5m}{2}+\sum_{j=1}^{m}p_j
 \overset{\eqref{eq:xor-completion-counts}}{\le}
       \frac{5m}{2}+\OPT(F).
\]
For the reverse bound, take an assignment satisfying $\OPT(F)$
equations. Select its truth blocks, of total size $\frac{5m}{2}$
by~\eqref{eq:xor-truth-budget}, and, for each
satisfied equation, its unique witness with that assignment. Such a
witness has no neighbor in the selected truth blocks. Only one witness
is selected per equation, so no star edge has both endpoints selected.
There are no edges between distinct equation gadgets. The resulting
independent set has size $\frac{5m}{2}+\OPT(F)$.

The matching bounds and the graph order $9m$ give
\begin{equation}\label{eq:xor-independent-set-optimum}
 \alphaI(G_F)=\frac{5m}{2}+\OPT(F),\qquad \vc(G_F)
 \overset{\eqref{eq:cover-independent-complement}}{=}9m-\alphaI(G_F)
 \overset{\eqref{eq:xor-independent-set-optimum}}{=}
 \frac{13m}{2}-\OPT(F).
\end{equation}
\end{proof}

\subsubsection{The anchor graph}
\label{subsubsec:anchor-graph}

Let $G$ be tripartite with specified independent parts $A,B,C$, and
let $\ell\ge1$ be an integer. Add pairwise disjoint sets
$K_A,K_B,K_C$ of size $\ell$ and distinct vertices $z_A,z_B,z_C$,
with all added vertices distinct and outside $V(G)$.
The graph $H_\ell(G)$ has vertex set
$V(G)\cup K_A\cup K_B\cup K_C\cup\{z_A,z_B,z_C\}$ and edge set
\[
 E(H_\ell(G))=E(G)\cup
 \bigcup_{X\in\{A,B,C\}}
 \left(\binom{K_X}{2}\cup
       \{\{u,v\}:u\in K_X,\ v\in X\cup\{z_X\}\}\right).
\]
Thus $K_X$ is a clique joined completely to $X\cup\{z_X\}$, and
there are no other new edges. We call $K_X$ the \emph{anchor clique}
of part $X$. This is the anchor graph of~\cite[Lemma~6]{BNS} and has
$|V(G)|+3\ell+3$ vertices.

\begin{corollary}[Anchors as long as the largest part]\label{cor:max-part-anchor}
If $G$ is tripartite with nonempty specified independent parts $A,B,C$ and
$\ell\ge\max\{|A|,|B|,|C|\}$, then
$\td(H_\ell(G))=\ell+\vc(G)+1$.
\end{corollary}

We derive this from the anchor construction in Appendix~\ref{app:anchor-proof},
including the short-anchor analysis.
Bonnet, Neuen, and Soko\l owski require $\ell\geq\vc(G)$ for their
identity~\cite[Lemma~6]{BNS}. Our corollary only requires $\ell$ to be at
least the largest part size, so it applies with $\ell=3m$ even though
$\vc(G_F)\overset{\text{Lem.~\ref{lem:xor-cover}}}{\geq}11m/2$.

\subsubsection{The treedepth identity and the hardness gap}
\label{subsec:np-td-gap}

\begin{proposition}\label{prop:xor-treedepth-identity}
For every normalized source system $F$, the graph $H=H_{3m}(G_F)$
has $18m+3$ vertices and satisfies
\begin{equation}\label{eq:xor-treedepth-identity}
 \td(H)=\frac{19m}{2}-\OPT(F)+1.
\end{equation}
\end{proposition}

\begin{proof}
Each source part has size $3m$ by~\eqref{eq:xor-equal-parts}, so
Corollary~\ref{cor:max-part-anchor} and Lemma~\ref{lem:xor-cover} give
\[
 \td(H)
 \overset{\text{Cor.~\ref{cor:max-part-anchor}}}{=}3m+\vc(G_F)+1
 \overset{\text{Lem.~\ref{lem:xor-cover}}}{=}
 \frac{19m}{2}-\OPT(F)+1.
\]
The graph has $9m+9m+3=18m+3$ vertices and $O(m^2)$ edges.
Repeated equations have distinct witness gadgets, so no parallel
edges arise. The source normalization changes size by a constant
factor, giving a polynomial construction.
\end{proof}

The treedepth reduction follows the chain
\[
 F\xrightarrow{\text{witness graph}}G_F
 \xrightarrow{\text{add anchors}}H_{3m}(G_F).
\]
Lemma~\ref{lem:xor-cover} encodes $\OPT(F)$ in the vertex-cover
number, and Proposition~\ref{prop:xor-treedepth-identity} converts
that number into treedepth. The source gap then gives
the two bounds below; Corollary~\ref{cor:td-bounded-degree} later
transfers them to bounded-degree graphs.

\begin{proof}[Dense treedepth gap: proof sketch]
Proposition~\ref{prop:xor-treedepth-identity} and the parity promise
give YES height at most $(17/2+\epsilon)m+1$ and NO height at least
$(9-\epsilon)m+1$. Their leading coefficient ratio tends to
$18/17$ as $\epsilon$ decreases. For any fixed $r<18/17$, choose
$\epsilon$ small enough and repeat the source a fixed number of
times so that the coefficient gap dominates the additive one.
Appendix~\ref{app:dense-np-td} gives the explicit error and
replication choices and checks the strict gap.
\end{proof}
